\section{LangFuse：本地部署的 API 监控工具}

\subsection{LangFuse 简介}

LangFuse 是一个可本地部署的 LLM API 监控工具，用于 trace 所有的模型调用过程。

\subsection{部署与配置}

\begin{enumerate}
    \item 下载 Docker Compose 文件并启动服务
    \item 本地访问 \texttt{localhost:3000}
    \item 创建组织和项目
    \item 获取 Public Key、Secret Key 和 Base URL
    \item 在 \texttt{.env} 中配置以上信息
\end{enumerate}

LightRAG \textbf{原生集成}了 LangFuse——只要检测到环境变量，就会自动将 OpenAI Client 替换为 LangFuse Client，在后台记录完整的 traces。

\subsection{通过 Traces 理解 Workflow}

在 LangFuse 的 Traces 面板中，可以看到 LightRAG 的所有 API 调用：

\textbf{Insert（Indexing）阶段}：
\begin{enumerate}
    \item 对文档创建 Embedding
    \item 第一次 Generation：提取实体和关系（Entity Extraction System/User Prompt）
    \item 第二次 Generation：Self Correction（继续提取遗漏的或格式不正确的内容）
    \item 大量 Embedding 调用：对节点和边创建 Embedding
\end{enumerate}

\textbf{Query（Retrieval）阶段}：
\begin{enumerate}
    \item Generation：提取 High-Level 和 Low-Level 关键词（Keywords Extraction Prompt）
    \item Embedding Search：基于关键词匹配实体和关系
    \item 图遍历：找到相关实体的邻居节点
    \item Generation：基于 Context + Query 生成最终回答（RAG Response Prompt）
\end{enumerate}

\begin{warningbox}{Response Cache}
LightRAG 默认启用 Response Cache——如果重复执行相同的 Query，不会再次调用 API。这在开发调试时需要注意。
\end{warningbox}

\subsection{本章小结}

LangFuse 提供了强大的可视化 trace 能力，是理解 Agent 项目内部工作机制的利器。

